\section{Value functions, local corrections and convex recourse: deferred
proofs}\label{app:recourse-convex}

This appendix contains the proofs deferred from
Sections~\ref{sec:valuefn}--\ref{sec:convexrecourse} and supplementary
results on certified convex recourse. The notation is that of these
sections. For a rational vector $w$ we write $\langle w\rangle$ for its
encoding length.

\subsection{Proof of Theorem~\ref{thm:valuefn}(b),(c)}

(b) Let $P_L(I)\ge I$ be the polynomial bound on the encoding lengths of
the $L_i^V$, and $P$ the polynomial such that a factor of $F_0$ or the oracle
at a rational point $w$ takes at most $P(I+\langle w\rangle)$ bit operations;
its output then has encoding length at most $P(I+\langle w\rangle)$. As
noted after Definition~\ref{def:growth}, point growth of $V$ with constant
$g$ gives weighted growth of $V$ with $\bar\kappa\le\kappa_V$. As in the
proof of Theorem~\ref{thm:approx}(b), CT succeeds at a trial with
$2^\mu\le6\sqrt{\kappa_V}$, and the number of table operations, factor
evaluations and oracle calls over all trials is at most
$c(J+1)pI^2(60p\sqrt{\kappa_V})^p$, now with $\abs{\mathcal A}+r\le2I$
factors. The exponents $e_i$ and $E$ of Section~\ref{sec:growth} are now
$O(P_L(I))$, so $J=O(P_L(I)+q)$, and the bit analysis of
Appendix~\ref{app:growth}, with $O(I)$ replaced by $O(P_L(I))$, shows that
every grid node of a trial with grading $2^{-\mu}$ has encoding length at
most $c(1+\mu2^\mu)(P_L(I)+q+1)$; moreover $1+\mu2^\mu\le c'\kappa_V$ for
the trials that run. Only the cost of one operation must be bounded anew,
because values of value factors have no common denominator. A table entry
is a sum of at most $\abs{\mathcal A}+r+n\le3I$ exact values and corrections
at one grid point $v$, each of encoding length at most
$P(I+\langle v\rangle)$, and $\langle v\rangle\le
cc'n\kappa_V(P_L(I)+q+1)$. So each addition, comparison, factor evaluation
and oracle call costs at most $(c''\kappa_V(I+q+1))^{C'}$ bit operations,
where $C'$ and $c''$ depend only on $P$ and $P_L$. Multiplying by the
number of operations, and using $I+q+1\ge2$ to absorb constant factors into
the exponent, gives the bound $f(p,\kappa_V)\kappa_V^C(I+q+1)^C$.

(c) Each run of CT on $V$ halts, because the proof of
Theorem~\ref{thm:approx}(a) uses only Definition~\ref{def:curvature}. Its
incumbent lifts to a feasible point $\hat x=(\hat v,y(\hat v))$
of~\eqref{eq:model} with $F(\hat x)-\beta\le2^{-q}$ and a certified
$\beta\le\OPT$. REC and the test of Proposition~\ref{prop:accept} are
applied to $F$ at $\hat x$, so validity and termination follow as in the
proof of Theorem~\ref{thm:exact}. For the bound, let $F$ have point growth
with constant $g$. If $L^V=0$, CT is exact at its first stage, $\hat x$ is
optimal, and REC returns an optimal point, which passes the test
(Lemma~\ref{lem:snap}). Otherwise point growth of $F$ is set growth with
$\mathcal S=\{(v^*,y^*)\}$, and Theorem~\ref{thm:transfer}(a) applies once
$2^{-q}\le\varepsilon_{\mathcal S}$, that is, for some
$q\le\poly(I)+O(\log\kappa_V)$, because $g\ge L^V/\kappa_V$ and
$\log(1/L^V)\le\poly(I)$. With (b) and $O(\log q)$ runs, absorbing powers of
$1+\log\kappa_V$ into $\kappa_V^{C_1}$ proves the bound.
\qed

\subsection{Supplements to Section~\ref{sec:local}}

\begin{proof}[Proof of Theorem~\ref{thm:cr-filter}]
(a) By~\eqref{eq:cr-oracle} and Lemma~\ref{lem:cr-cell} applied to $x^*$,
\[
\min_{\text{corners}}\underline V-e\le\min_{\text{corners}}V_B-e_B(C)\le\OPT\le U.
\]

(b) Lemma~\ref{lem:cr-cell} at $x^*$ gives a corner $v^\circ$ of the
optimizer's cell with $V_B(v^\circ)\le\OPT+e$. Hence
$U\le F(x^{v^\circ})\le\underline V(v^\circ)+\varepsilon_{\mathrm{or}}\le V_B(v^\circ)+\varepsilon_{\mathrm{or}}\le\OPT+e+\varepsilon_{\mathrm{or}}$.
A retained cell has a corner $v$ with $\underline V(v)\le U+e$, and
$V_B(v)\le F(x^v)\le\underline V(v)+\varepsilon_{\mathrm{or}}\le U+e+\varepsilon_{\mathrm{or}}\le\OPT+2e+2\varepsilon_{\mathrm{or}}$.

(c) For $x_B\in C$, Lemma~\ref{lem:cr-cell} and $\underline V\le V_B$ give
$F(x)\ge\min_{\text{corners}}V_B-e\ge\min_{\text{corners}}\underline V-e\ge\lambda$.
Every corner satisfies $\underline V(v)\ge F(x^v)-\varepsilon_{\mathrm{or}}\ge U-\varepsilon_{\mathrm{or}}$, so
$\lambda\ge U-e-\varepsilon_{\mathrm{or}}$.
\end{proof}

\begin{proposition}[Only the error oscillation matters]\label{prop:cr-osc}
Let $\mathcal Q$, $e$ and $\varepsilon_{\mathrm{or}}$ be as in Theorem~\ref{thm:cr-filter}, and
suppose some cell of $\mathcal Q$ contains $x^*_B$ for a global optimizer
$x^*$. Let $\mathcal V$ be the set of all corners and
$\tilde V:\mathcal V\to\R$ a function such that, for some constants
$\delta_-\le\delta_+$ with $\delta_+-\delta_-\le\varepsilon_{\mathrm{or}}$,
\[
\delta_-\le V_B(v)-\tilde V(v)\le\delta_+\qquad(v\in\mathcal V).
\]
The constants need not be known. Retain $C$ if and only if
$\min_{\text{corners of }C}\tilde V\le\min_{\mathcal V}\tilde V+e+\varepsilon_{\mathrm{or}}$.
Then every cell containing $x^*_B$ for a global optimizer $x^*$ is
retained, and every retained cell has a corner with
$V_B(v)\le\OPT+2e+2\varepsilon_{\mathrm{or}}$.
\end{proposition}

\begin{proof}
Let $C^*$ contain $x^*_B$. By Lemma~\ref{lem:cr-cell},
\[
\min_{\text{corners of }C^*}\tilde V\le\min_{\text{corners of }C^*}V_B-\delta_-
\le\OPT+e-\delta_-,
\]
while $V_B\ge\OPT$ gives $\min_{\mathcal V}\tilde V\ge\OPT-\delta_+$. The
difference is at most $e+\delta_+-\delta_-\le e+\varepsilon_{\mathrm{or}}$, so $C^*$ is retained.
If $C$ is retained, it has a corner $v$ with
$V_B(v)\le\tilde V(v)+\delta_+\le\min_{\mathcal V}\tilde V+e+\varepsilon_{\mathrm{or}}+\delta_+$, and
$\min_{\mathcal V}\tilde V\le\min_{\text{corners of }C^*}\tilde V\le\OPT+e-\delta_-$.
Thus $V_B(v)\le\OPT+2e+\varepsilon_{\mathrm{or}}+(\delta_+-\delta_-)\le\OPT+2e+2\varepsilon_{\mathrm{or}}$.
\end{proof}

\begin{proof}[Proof of Proposition~\ref{prop:star}]
(A) Write $\omega_A=mh^2/4-1-\epsilon>0$. All coordinate second derivatives
equal $2$ and $F_A$ is quadratic, so $L=2$ is valid. For fixed $x$, each leaf
term $y^2-hxy$ is strictly convex with minimizer $hx/2\in[0,1]$ and minimum
$-h^2x^2/4$. Hence $\min_yF_A(x,y)=(1-mh^2/4)x^2+\omega_Ax=:V(x)$, which is
strictly concave because $mh^2>4$, with $V(0)=0$ and $V(1)=-\epsilon$. For
$x\in(0,1)$, strict concavity gives $V(x)>(1-x)V(0)+xV(1)\ge-\epsilon$. So
$x=1$ is the unique minimizing center and $y_i=h/2$ the unique leaves;
$\OPT=-\epsilon$. The origin has $F_A=0=\OPT+\epsilon$ and squared distance
$1+mh^2/4$ from the minimizer, which bounds $g$.

On the grid, $y(y-hx)\ge0$ for $y\in G$ and $x\in[0,1]$, since either $y=0$
or $y\ge h\ge hx$. Also $x^2+\omega_Ax\ge0$ for $x\ge0$. Hence $F_A\ge0$ on
$G^{1+m}$ with equality at the origin, so $U_G=0$, and every leaf's grid
minimum is $0$. The point $(1,h/2)$ has $x$ at the right endpoint and $h/2$
interior to $[0,h]$, so $C$ is the only grid cell of the bag containing it.
At its corners, $m^F(x,y_1)=x^2+\omega_Ax+(y_1^2-hxy_1)$. For $x=1-h$ this
equals $(1-h)(1-h+\omega_A)=(1-h)(mh^2/4-h-\epsilon)$ at $y_1=0$ and exceeds
it by $h^3$ at $y_1=h$. For $x=1$ it equals $mh^2/4-\epsilon$ at both
$y_1=0,h$, which exceeds the previous value by
$h(mh^2/4-\epsilon)+h(1-h)>0$. This proves the corner minimum; it also equals
$\min\{V^G(1-h),V^G(1)\}$, because every leaf's grid minimum is attained at
$0$. Since $U_G=0$ and $L=2$, the threshold for $c$ is as stated, and
$(1+m)h^2/4\big/\bigl[(1-h)(mh^2/4-h-\epsilon)\bigr]\to(1-h)^{-1}$.

(B) The Hessian $H$ has $H_{xx}=2+2m/d^2=4$, $H_{y_iy_i}=2$,
$H_{xy_i}=-2/d$, and zero elsewhere; thus $L=4$ is valid. Vectors with zero
$x$-component and leaf components summing to zero are eigenvectors with
eigenvalue $2$. On the span of $e_x$ and $m^{-1/2}\mathbf 1_y$, $H$ acts as
$\bigl(\begin{smallmatrix}4&-2\\-2&2\end{smallmatrix}\bigr)$, with
eigenvalues $3\pm\sqrt5$. Since $F_B\ge0$ vanishes only at
$z^*\in(0,1)^{1+m}$ and $F_B(z)=\frac12(z-z^*)^{\T}H(z-z^*)$, growth holds
with $g=(3-\sqrt5)/2$ and $\kappa=L/g=2(3+\sqrt5)$.

Write $t=x-1/2\in h\Z\cap[-1/2,1/2]$ for the center grid values (here
$1/2\in G$ because $j\ge1$). A leaf's ideal value $\iota(t)=h/2+t/d$ lies in
$[0,h]$ because $|t|/d\le1/(2d)\le h/4$. Its distance to $G$ is therefore
$h/2-|t|/d$, and the grid value function of the center is
\[
V^G(t)=t^2+m\Bigl(\frac h2-\frac{|t|}d\Bigr)^2=2t^2-dh|t|+\frac{d^2h^2}4 .
\]
On the grid $|t|\le1/2\le dh/4$, and $\tau\mapsto2\tau^2-dh\tau$ decreases on
$[0,dh/4]$, so $U_G=V^G(\pm1/2)=\frac12-\frac{dh}2+\frac{d^2h^2}4$. The
optimizer's bag projection $(1/2,h/2)$ lies on the common edge of the cells
$[1/2-h,1/2]\times[0,h]$ and $[1/2,1/2+h]\times[0,h]$ and in no other cell.
At $t=0$ both leaf nodes $0,h$ are at distance $h/2$ from $\iota$, so
$m^F=d^2h^2/4$. At $t=\pm h$ the nearer node gives
$m^F=V^G(h)=2h^2-dh^2+d^2h^2/4$ and the farther adds $2h^2/d$. Since $d\ge2$,
each cell has corner minimum $d^2h^2/4-(d-2)h^2=\min\{V^G(0),V^G(\pm h)\}$,
and subtracting $U_G$ gives the stated difference. For $h\le1/4$ this
difference is at least $dh/4+2h^2-1/2$, which grows linearly in $d=\sqrt m$.
Discarding both center intervals adjacent to $x=1/2$ removes the optimizer.
\end{proof}

\begin{example}[An instance with $h=\frac12$]\label{ex:cr-star32}
Take $h=\frac12$, $m=32$, $\epsilon=\frac1{16}$ in
Proposition~\ref{prop:star}(A):
$F=x^2-\frac x2\sum y_i+\sum y_i^2+\frac{15}{16}x$ on $[0,1]^{33}$.
The exact leaf recourse is $y_i=x/4$, with $V(x)=-x^2+\frac{15}{16}x$, so
the unique optimizer is $x=1$, $y_i=\frac14$, value $-\frac1{16}$. On
$G=\{0,\frac12,1\}$ the grid values $V^G$ at $x=0,\frac12,1$ are
$0,\frac{23}{32},\frac{31}{16}$, and $U_G=0$. The four corner
min-marginals of $C=[\frac12,1]\times[0,\frac12]$ are
$\frac{23}{32},\frac{27}{32},\frac{31}{16},\frac{31}{16}$. The bag-local
budget $\frac18$ rejects $C$; the global correction $\frac{33}{16}$ keeps it.
The grid error $V^G(x)-V(x)=m\,\dist(x/4,G)^2$ varies from $0$ at $x=0$ to
$m/16$ at $x=1$; no normalization of messages removes this variation.
\end{example}

\subsection{Certificates: energy, existence, evaluation and height}

In the setting of Lemma~\ref{lem:leaf}, put
$\mathcal W=\{d:G_sd=0\ (s\in\mathcal E)\}$. For every coordinate $j$,
\begin{equation}\label{eq:cv-energy}
(B^{\T}CB)_{jj}
=-\min_{d\in\mathcal W}\bigl(d^{\T}Cd+2d^{\T}\Gamma e_j\bigr).
\end{equation}
Indeed, the function $d\mapsto d^{\T}Cd+2d^{\T}\Gamma e_j$ is convex. By
Lemma~\ref{lem:leaf}(b), $Be_j\in\mathcal W$, and by the $w$-linear part
$CB+\Gamma+G^{\T}\Lambda=0$ of~\eqref{eq:leaf-stat}, for $d\in\mathcal W$
half the gradient at $Be_j$ satisfies
$(CBe_j+\Gamma e_j)^{\T}d=-\sum_{s\in\mathcal E}(\Lambda_se_j)(G_sd)=0$. A
convex function whose gradient at a point of a linear subspace is orthogonal
to that subspace attains its minimum over the subspace there, so the minimum
equals $(B^{\T}CB)_{jj}+2(B^{\T}\Gamma)_{jj}=-(B^{\T}CB)_{jj}$.
Formula~\eqref{eq:cv-energy} says that a piece cancels direct curvature
exactly by the energy that the private response spends to follow a unit
change of $w_j$, using only directions that keep every constraint with a
nonzero multiplier active. If all private coordinates coupled to $w_j$ are
held at bounds with nonzero multipliers, the piece cancels nothing.


\begin{proof}[Proof of Proposition~\ref{prop:vf-curv}(b)]
Let $J=[t_-,t_+]$ with $t_-<t_+$; a point interval is trivial. The sets
$J_{\Pi'}=\{t\in J:w+te_j\in\Pi'\}$ are closed intervals that cover $J$. The
union of the nondegenerate ones is closed, and its complement in $J$ is
relatively open and contained in the finitely many degenerate ones, which
are points; hence it is empty. The endpoints of the nondegenerate
$J_{\Pi'}$ therefore split $J$ into finitely many closed subintervals. Some
nondegenerate $J_{\Pi'}$ contains the midpoint of a given subinterval, and
then contains the whole subinterval, because no endpoint of $J_{\Pi'}$ lies
strictly inside it. On that subinterval $\psi(w+te_j)=q_{\Pi'}(w+te_j)$ has
second derivative $-(B_{\Pi'}^{\T}CB_{\Pi'})_{jj}\le-\sigma_j$ by
Lemma~\ref{lem:leaf}, so $\chi(t)=\psi(w+te_j)+\sigma_jt^2/2$ is concave
there. The function $\chi$ is continuous, and at each breakpoint the jump
$\chi'_+-\chi'_-$ equals that of the concave function $t\mapsto\psi(w+te_j)$,
which is nonpositive. So $\chi'_+$ is nonincreasing, and $\chi$ is concave.
\end{proof}

\begin{proposition}[Finite certificates exist]\label{prop:cert-exist}
Let $Y=[\ell^Y,u^Y]$ with $\ell^Y<u^Y$, and let $\Pi\subseteq\R^k$ be a
rational box with nonempty interior. Then $\psi$ has a certificate on $\Pi$
in the sense of Definition~\ref{def:leaf}.
\end{proposition}

\begin{proof}
A \emph{pattern} $J=(J_\ell,J_f,J_u)$ is a partition of the private
coordinates. Let $\mathcal M_J$ be the set of
$(w,y,\lambda^{\mathrm{low}},\lambda^{\mathrm{up}})$ with $w\in\Pi$,
$\ell^Y\le y\le u^Y$, $\lambda^{\mathrm{low}},\lambda^{\mathrm{up}}\ge0$,
$Cy+\Gamma w+c=\lambda^{\mathrm{low}}-\lambda^{\mathrm{up}}$, and
$y_i=\ell^Y_i$, $\lambda^{\mathrm{up}}_i=0$ for $i\in J_\ell$;
$\lambda^{\mathrm{low}}_i=\lambda^{\mathrm{up}}_i=0$ for $i\in J_f$;
$y_i=u^Y_i$, $\lambda^{\mathrm{low}}_i=0$ for $i\in J_u$. Each
$\mathcal M_J$ is a rational polyhedron, and every point of it satisfies the
KKT conditions, including complementarity. Its projection $\Pi_J$ to the $w$
coordinates is a rational polytope. Every $w\in\Pi$ has an optimal $y$ and,
since the constraints are linear, KKT multipliers; these fit some pattern.
Thus the $\Pi_J$ cover $\Pi$. The union of the full-dimensional $\Pi_J$ is
closed and contains $\Pi$ minus finitely many lower-dimensional polytopes, a
dense subset of $\Pi$; hence it equals $\Pi$.

Triangulate each full-dimensional $\Pi_J$ into $k$-simplices with rational
vertices. At each vertex $w'$ of a simplex $\Sigma$, choose a rational point
$(w',y_{w'},\lambda_{w'})\in\mathcal M_J$. On $\Sigma$ define
$(\bar y,\lambda)$ by affine interpolation of these vertex values in
barycentric coordinates. For $w\in\Sigma$ the point
$(w,\bar y(w),\lambda(w))$ is a convex combination of points of the convex
set $\mathcal M_J$, so it lies in $\mathcal M_J$. Hence $(\bar y,\lambda)$ is
a certificate on $\Sigma$.

Let $\mathcal H$ be the finite set of hyperplanes spanned by facets of all
these simplices. Starting from $\Pi$, split every current leaf whose interior
meets a hyperplane of $\mathcal H$ by that hyperplane, until no leaf interior
meets any of them. Both halves of every split have nonempty interior, and no
hyperplane splits two nodes on one root path, so the process stops. Let
$\Pi'$ be a final leaf. Its interior is connected and misses every hyperplane
of $\mathcal H$, so for each simplex $\Sigma$ either
$\operatorname{int}\Pi'\subseteq\operatorname{int}\Sigma$ or
$\operatorname{int}\Pi'\cap\Sigma=\emptyset$. The full-dimensional simplices
cover $\Pi$, so the first case occurs for some $\Sigma$, and
$\Pi'\subseteq\Sigma$. Restricting the certificate of $\Sigma$ to $\Pi'$
completes the proof.
\end{proof}

The same argument applies to a polytope $Y$, with patterns given by sets of
active rows.

\begin{lemma}[Exact tables with value factors]\label{lem:cv-bits}
For rational $w\in\R^{k_t}$, the value $\psi_t(w)$ and a minimizer
$y^{(t)}(w)$ are rationals of encoding length $\poly(I+\langle w\rangle)$,
computable in that time, where $I$ is the encoding length of the data. In
CT applied to $V$, every message, bag-table entry and min-marginal is a sum
of at most $m_0=\abs{\mathcal A}+r+\abs{\mathcal K}$ factor values,
value-factor values and corrections at one grid point $v$, hence has
encoding length at most $m_0\poly(I+\langle v\rangle)$.
\end{lemma}

\begin{proof}
A convex quadratic program with rational data is solvable exactly in
polynomial time, and it has a rational optimal solution and value of
polynomial size \cite{KozlovTarasovKhachiyan1980}. If a certificate is
supplied, the value is instead $q_{\Pi'}(w)$ and the minimizer
$\bar y(w)$ for a leaf $\Pi'$ containing $w$, found by descending the
partition tree. In the message recurrence and in the bag tables, every entry
equals the value of a sum of factor values and unary corrections at one
assignment, with each factor and each correction appearing at most once. If
each of $m_0$ rationals has a numerator of absolute value at most $2^{b'}$
and a denominator at most $2^{b'}$, their sum has denominator at most
$2^{m_0b'}$ and numerator of absolute value at most $m_02^{m_0b'}$.
Comparisons and additions on such numbers are polynomial in their length. No
common denominator across different active sets is needed.
\end{proof}

\begin{lemma}[Height of an optimizer over a polytope]\label{lem:cv-height}
Let $\mathcal P=\{x\in\R^{n'}:Ax\le a\}$ be a nonempty bounded polyhedron
with integral $A$ and $a$, and let
$F_{\mathcal P}(x)=(\frac12x^{\T}Hx+c^{\T}x+c_0)/\Delta_{\mathcal P}$ with
integral symmetric $H$, integral $c$ and $c_0$, and a positive integer
$\Delta_{\mathcal P}$. Put
$\beta_H=\max\{1,\max_{ij}|H_{ij}|,\max_{sj}|A_{sj}|\}$ and
$\Lambda_0=(2n'\beta_H^2)^{n'}$. Then some global minimizer $x^*$ of
$F_{\mathcal P}$ on $\mathcal P$ has $x^*=z/\rho$ with $z\in\Z^{n'}$ and an
integer $1\le\rho\le\Lambda_0$; if the minimizer is unique, it has this
form. The minimum $\min_{\mathcal P}F_{\mathcal P}$ is a rational with
denominator at most $2\Delta_{\mathcal P}\Lambda_0^2$.
\end{lemma}

\begin{proof}
Global minimizers exist because $\mathcal P$ is compact. Among them choose
$x^*$ whose active set $\mathcal E=\{s:A_sx^*=a_s\}$ has maximal
cardinality, and put $\mathcal W=\{d:A_{\mathcal E}d=0\}$. For
$d\in\mathcal W$ and small $\epsilon>0$ both $x^*\pm\epsilon d$ lie in
$\mathcal P$, so optimality gives $\nabla F_{\mathcal P}(x^*)^{\T}d=0$ and
$d^{\T}Hd\ge0$.

Suppose some $d\in\mathcal W\setminus\{0\}$ has $Hd\perp\mathcal W$. Then
$d^{\T}Hd=0$, so $F_{\mathcal P}(x^*+td)=F_{\mathcal P}(x^*)$ for all real
$t$. As $\mathcal P$ is bounded, $\bar t=\max\{t\ge0:x^*+td\in\mathcal P\}$
is finite. The rows of $\mathcal E$ stay active along the line. For every
$t>\bar t$ some row is violated at $x^*+td$; as there are finitely many
rows, one row $s\notin\mathcal E$ is violated for $t$ arbitrarily close to
$\bar t$, and continuity and feasibility give $A_s(x^*+\bar td)=a_s$. So
$x^*+\bar td$ is a global minimizer with a strictly larger active set, a
contradiction. Hence no nonzero $d\in\mathcal W$ has $Hd\perp\mathcal W$.
Since $H$ is positive semidefinite on $\mathcal W$, it is then positive
definite on $\mathcal W$: a nonzero $d\in\mathcal W$ with $d^{\T}Hd=0$ would
minimize the convex form $d\mapsto d^{\T}Hd$ over $\mathcal W$, so
$Hd\perp\mathcal W$.

Choose a maximal linearly independent set $\mathcal E'\subseteq\mathcal E$
of rows, so that $\mathcal W=\{d:A_{\mathcal E'}d=0\}$. Since
$\nabla F_{\mathcal P}(x^*)\perp\mathcal W$, the vector $Hx^*+c$ lies in the
row space of $A_{\mathcal E'}$: there is $\mu$ with
$Hx^*+c+A_{\mathcal E'}^{\T}\mu=0$, and $A_{\mathcal E'}x^*=a_{\mathcal E'}$.
The saddle matrix
$S=\left(\begin{smallmatrix}H&A_{\mathcal E'}^{\T}\\A_{\mathcal E'}&0\end{smallmatrix}\right)$
is nonsingular: if $Hx+A_{\mathcal E'}^{\T}\mu'=0$ and $A_{\mathcal E'}x=0$,
then $x\in\mathcal W$ and $x^{\T}Hx=-(A_{\mathcal E'}x)^{\T}\mu'=0$, so $x=0$,
and then $\mu'=0$ because the rows of $A_{\mathcal E'}$ are independent. The
order of $S$ is at most $2n'$, its entries are integers of absolute value at
most $\beta_H$, and so each row $S_i$ has norm at most $\sqrt{2n'}\,\beta_H$.
Hadamard's inequality for rows, $\abs{\det S}\le\prod_i\norm{S_i}$, gives
$\abs{\det S}\le(\sqrt{2n'}\,\beta_H)^{2n'}=\Lambda_0$. Cramer's rule applied
to $S(x^*,\mu)=(-c,a_{\mathcal E'})$, whose right side is integral, gives
$x^*=z/\rho$ with $\rho=\abs{\det S}$. Finally
$F_{\mathcal P}(x^*)=(z^{\T}Hz+2\rho c^{\T}z+2\rho^2c_0)/(2\rho^2\Delta_{\mathcal P})$.
If the minimizer is unique, it is $x^*$.
\end{proof}

\begin{proof}[Proof of the exact-output part of Theorem~\ref{thm:cv}(iii)]
Fix the integer retained coordinates at any integer values. The remaining
problem minimizes a quadratic over the bounded rational polytope formed by
the continuous retained box and the sets $Y_t$. Scale its constraint rows
to integers and multiply the objective by a common denominator
$\Delta_{\mathcal P}$ of all coefficients of $F$. The fixed integer values
change only the linear and constant terms, which stay integral, so the
numbers $\beta_H$ and $\Delta_{\mathcal P}$ of Lemma~\ref{lem:cv-height} are
the same for every integer slice. That lemma gives computable integers
$\Lambda_0$ and $\Omega_0=2\Delta_{\mathcal P}\Lambda_0^2$ of
polynomial encoding length: every slice optimum, in particular $\OPT$, has
denominator at most $\Omega_0$, and a unique minimizer has continuous
coordinates with denominators at most $\Lambda_0$.

Run CT on $V$ with $\varepsilon=2^{-q}$ for $q=1,2,4,\ldots$. After each run,
with incumbent $\hat v$ and bound $\beta$, replace each continuous coordinate
of $\hat v$ by the unique rational with denominator at most $\Lambda_0$
within distance $1/(4\Lambda_0^2)$, if one exists; continued fractions find
it in polynomial time. If every such coordinate is replaced and the result
$\tilde v$ lies in $X_{\mathcal K}$, compute $y(\tilde v)$ and
$V(\tilde v)=a/W$ in lowest terms, and stop with $(\tilde v,y(\tilde v))$ if
$V(\tilde v)-\beta<1/(\Omega_0W)$. By Proposition~\ref{prop:accept} with the
denominator bound $\Omega_0$ for $\OPT$, a returned point is optimal: if
$V(\tilde v)>\OPT$, then $V(\tilde v)-\beta\ge V(\tilde v)-\OPT\ge1/(\Omega_0W)$.
This uses neither growth nor uniqueness.

If $F$ has point growth with constant $g$, its minimizer $(v^*,y^*)$ is
unique, so the continuous coordinates of $v^*$ have denominators at most
$\Lambda_0$. Once
$\varepsilon\le\min\{1/(2\Omega_0^2),\,g/(32\Lambda_0^4)\}$,
Lemma~\ref{lem:valuefunction}(b) gives
$\norm{\hat v-v^*}^2\le(V(\hat v)-\OPT)/g\le1/(32\Lambda_0^4)$. Every
coordinate of $\hat v$ is then within $1/(4\Lambda_0^2)<1$ of the
corresponding coordinate of $v^*$, so integer coordinates agree, and two
distinct rationals with denominators at most $\Lambda_0$ differ by at least
$1/\Lambda_0^2$, so $\tilde v=v^*$. Then $V(\tilde v)=\OPT=a/W$ with
$W\le\Omega_0$ and $V(\tilde v)-\beta\le\varepsilon<1/(\Omega_0W)$, so the
procedure stops. Since $1/g\le\kappa_V/L^+$ and $\log(1/L^+)$,
$\log\Lambda_0$, $\log\Omega_0$ are polynomial in $I$, this happens with
$q\le\poly(I)+O(\log\kappa_V)$, and doubling overshoots by a factor at most
two. Each run costs at most $f(p,\kappa_V)\kappa_V^C(I+q+1)^C$ bit operations by
the approximation part, the reconstruction and the evaluation of
$V(\tilde v)$ cost $\poly(I+q)$, and there are $O(\log q)$ runs; absorbing
the powers of $1+\log\kappa_V$ into $\kappa_V^{C_1}$ gives the bound
$f(p,\kappa_V)\kappa_V^{C_1}(I+1)^{C_1}$.
\end{proof}

\subsection{Proof of Theorem~\ref{thm:cv-recog}}

\emph{Common central gradient.} Let $\hat y$ be an optimizer at $w_0$. For
any optimizer $y'$ at $w_0$, the function
$t\mapsto\phi(w_0,\hat y+t(y'-\hat y))$ is a convex quadratic that is
constant on $[0,1]$, because the optimal set is convex. Its quadratic
coefficient $(y'-\hat y)^{\T}C(y'-\hat y)$ vanishes, so $C(y'-\hat y)=0$
since $C\succeq0$. All central optimizers therefore have the same gradient
$\zeta_0=C\hat y+\Gamma w_0+c$; this is a classical property of convex
programs \cite{Mangasarian1988}. Let $I_\ell=\{i:\zeta_{0i}>0\}$,
$I_u=\{i:\zeta_{0i}<0\}$ and $I_0=\{i:\zeta_{0i}=0\}$. By the box KKT
conditions, every central optimizer has $y_i=\ell^Y_i$ for $i\in I_\ell$ and
$y_i=u^Y_i$ for $i\in I_u$.

\emph{Necessity.} We use one fact: if an affine function is nonnegative on
$\Pi$ and vanishes at an interior point, it vanishes identically. Let $\bar y$
be an affine optimal selector and write $\zeta(w)=C\bar y(w)+\Gamma w+c$, an
affine map. Since $\bar y(w_0)$ is a central optimizer, $\zeta(w_0)=\zeta_0$.
For $i\in I_\ell$, $\bar y_i-\ell^Y_i\ge0$ on $\Pi$ and vanishes at $w_0$,
so $\bar y_i\equiv\ell^Y_i$; optimality at every $w$ then needs
$\zeta_i\ge0$ on $\Pi$. Symmetrically, $\bar y_i\equiv u^Y_i$ and
$\zeta_i\le0$ for $i\in I_u$. Let $j\in I_0$, so $\zeta_j(w_0)=0$. If
$\bar y_j\equiv\ell^Y_j$, optimality needs $\zeta_j\ge0$ on $\Pi$, so
$\zeta_j\equiv0$; the case $\bar y_j\equiv u^Y_j$ is symmetric. Otherwise
$\bar y_j$ is not identically at a bound, so by the fact above it attains no
bound at an interior point. Then optimality gives $\zeta_j=0$ on the interior
of $\Pi$, hence $\zeta_j\equiv0$. Thus every affine optimal selector
satisfies
\begin{equation}\label{eq:recog-system}
\begin{array}{ll}
\bar y_i\equiv\ell^Y_i,\ \zeta_i\ge0\text{ on }\Pi & (i\in I_\ell),\\
\bar y_i\equiv u^Y_i,\ \zeta_i\le0\text{ on }\Pi & (i\in I_u),\\
\zeta_j\equiv0,\ \ell^Y_j\le\bar y_j\le u^Y_j\text{ on }\Pi & (j\in I_0).
\end{array}
\end{equation}

\emph{Sufficiency.} If an affine $\bar y$ satisfies~\eqref{eq:recog-system},
define lower multipliers $\zeta_i$ on $I_\ell$ and upper multipliers
$-\zeta_i$ on $I_u$, and zero otherwise. These are affine and nonnegative on
$\Pi$, and satisfy stationarity and complementarity identically. This is a
certificate on $\Pi$; by Lemma~\ref{lem:leaf}(a), $\bar y(w)$ is optimal for
every $w\in\Pi$.

\emph{Linear program.} Write $\bar y(w)=\bar y_0+Bw$. Then
$\zeta(w)=(C\bar y_0+c)+(CB+\Gamma)w$ has coefficients linear in
$(\bar y_0,B)$. Identities in~\eqref{eq:recog-system} are linear equations.
Each condition ``an affine function is nonnegative on $\Pi$'' is replaced by
its Farkas certificate $\pi\ge0$, $A^{\T}\pi=-(\text{linear part})$,
$a^{\T}\pi\le(\text{constant part})$, which is linear in $(\bar y_0,B,\pi)$.
For a box $\Pi$, auxiliary variables bounding $|B_{ij}|$ and
$|(CB+\Gamma)_{ij}|$ give the exact ranges instead. The resulting system has
polynomial size. It is feasible exactly when an affine optimal selector
exists. Exact convex quadratic programming \cite{KozlovTarasovKhachiyan1980}
and exact linear programming \cite{GrotschelLovaszSchrijver1988} run in
polynomial time and return rational solutions of polynomial size. If $k=0$,
any central optimizer is a constant selector. \qed

\medskip
Fixing the active bounds of an arbitrary central optimizer is not enough
when $C$ is singular. For $(y_1+y_2-w)^2$ with $y\in[0,1]^2$ and
$w\in[0,2]$, the central optimizer $(0,1)$ would fix both coordinates, yet
the common central gradient is zero and the theorem returns the selector
$(w/2,w/2)$. For $(y-w)^2$ with $y\in[0,1]$ and $w\in[-1,2]$, the clipped
response is not affine, and the linear program is infeasible.

\subsection{Affine condensation and a negative-curvature regime}

\begin{lemma}[Growth transfer]\label{lem:cv-growth}
In the model of Definition~\ref{def:cv-model}, suppose
$F(v,y)-\OPT\ge g\bigl(\norm{v-v^*}^2+\norm{y-y^*}^2\bigr)$ on the feasible
set, with $g>0$. Then for every $v\in X_{\mathcal K}$ and every choice of
optimal responses $y(v)$,
\[
V(v)-\OPT\ge g\bigl(\norm{v-v^*}^2+\norm{y(v)-y^*}^2\bigr).
\]
If block $t$ has a one-leaf certificate $\bar y_t$ on $\Pi_t$, then
$y^{*(t)}=\bar y_t(v^*_{E_t})$.
\end{lemma}

\begin{proof}
Evaluate the growth inequality at $(v,y(v))$ and use $F(v,y(v))=V(v)$.
Replacing $y^{*(t)}$ by the optimal response $\bar y_t(v^*_{E_t})$ gives a
point of value $V(v^*)=\OPT$, which equals $(v^*,y^*)$ because growth makes
the minimizer unique.
\end{proof}

\begin{corollary}[Affine condensation]\label{cor:cv-affine}
Suppose every block has a one-leaf certificate $\bar y_t$ on $\Pi_t$, the
retained coordinates are continuous, and $F$ has point growth with
constant $g$. Let $\mathcal J$ be the matrix of the linear part of the affine
map $v\mapsto(v,\bar y(v))$, where $\bar y(v)$ collects the responses
$\bar y_t(v_{E_t})$, so that
$\norm{\mathcal Jd}^2=\norm d^2+\sum_t\norm{B_td_{E_t}}^2$. Let
$H_{\mathrm{red}}=\nabla^2V$ and $\nu=\max\{0,-\lambda_{\min}(\nabla^2F)\}$.
\begin{enumerate}
\item[(a)] $V(v)-\OPT\ge g\norm{\mathcal J(v-v^*)}^2$ and
$H_{\mathrm{red}}\succeq-\nu\mathcal J^{\T}\mathcal J$.
\item[(b)] If $\delta\in\R^{\mathcal K}$ is positive with
$\mathcal J^{\T}\mathcal J\succeq\diag(\delta)$, choose dyadic $\omega_i>0$
with $1\le\omega_i^2\delta_i<4$ and substitute $v=\diag(\omega)v'$. The new
box is rational, the interaction graph is unchanged, the growth constant in
$v'$ is at least $g$, and the upper coordinate curvature in $v'$ is at most
$4\max_i\max\{(H_{\mathrm{red}})_{ii},0\}/\delta_i$.
\end{enumerate}
\end{corollary}

\begin{proof}
(a) By Lemma~\ref{lem:cv-growth}, $y(v)=\bar y(v)$ and $y^*=\bar y(v^*)$,
so $\norm{v-v^*}^2+\norm{y(v)-y^*}^2=\norm{\mathcal J(v-v^*)}^2$. Since
$V(v)=F(v,\bar y(v))$ is a quadratic composed with an affine map,
$H_{\mathrm{red}}=\mathcal J^{\T}\nabla^2F\,\mathcal J\succeq-\nu\mathcal J^{\T}\mathcal J$;
by Lemma~\ref{lem:leaf}(b) it equals $H_0$ minus the matrices
$B_t^{\T}C_tB_t$ placed in the rows and columns $E_t$. (b) Growth in $v'$
is at least $g\sum_i\omega_i^2\delta_i(v'_i-v'^*_i)^2\ge g\norm{v'-v'^*}^2$,
and the diagonal entry in $v'$ is
$\omega_i^2(H_{\mathrm{red}})_{ii}<4(H_{\mathrm{red}})_{ii}/\delta_i$ when it
is positive.
\end{proof}

Rescaling an integer coordinate changes its lattice, so part~(b) is stated
for continuous coordinates. The inequality in (a) bounds negative curvature
only; it gives no upper bound on $(H_{\mathrm{red}})_{ii}$.

\begin{corollary}[A negative-curvature regime]\label{cor:cv-nu}
In Theorem~\ref{thm:cv}(iii), if the full Hessian of $F$ has
$\nu=\max\{0,-\lambda_{\min}\}>0$ and $L^+\le C_0\nu$, then
$\kappa_V\le\max\{1,C_0\}\max\{1,\nu/g\}$. Membership in this regime can be
tested up to a factor two: a rational $\hat\nu$ with $\nu\le\hat\nu<2\nu$ is
computable in polynomial time, and $L^+\le C_0\hat\nu/2$ implies
$L^+\le C_0\nu$.
\end{corollary}

\begin{proof}
The bound on $\kappa_V$ is immediate, and $\hat\nu/2<\nu$ gives the
implication. For $\hat\nu$, write $H=\nabla^2F$ and let $\Delta_H$ be a
positive integer with $\Delta_HH$ integral. First test $H\succeq0$ exactly;
if it fails, $\nu>0$. Start from $\hat\nu_0=\max_i\sum_j|H_{ij}|\ge\norm
H\ge\nu$ and halve $\hat\nu$ while $H+(\hat\nu/2)I\succeq0$. The final
$\hat\nu$ satisfies $H+\hat\nu I\succeq0$ and $H+(\hat\nu/2)I\not\succeq0$,
that is, $\nu\le\hat\nu<2\nu$. If $d=\operatorname{rank}H$, the product of
the nonzero eigenvalues is, up to sign, a coefficient of the characteristic
polynomial with denominator dividing $\Delta_H^d$, so its absolute value is
at least $\Delta_H^{-d}$. Each nonzero eigenvalue is at most $\hat\nu_0$ in
absolute value, so $\nu\ge\Delta_H^{-d}\hat\nu_0^{1-d}$, and the number of
halvings is $\poly(I)$. Each test is an exact symmetric elimination.
\end{proof}

The algorithm itself never uses $\nu$ or $\hat\nu$; it computes $L^+$. The
hypothesis $L^+\le C_0\nu$ is substantive: Proposition~\ref{prop:cv-limit}
shows that it can fail for every admissible choice of private blocks.

\subsection{Examples}

\begin{example}[Three pieces, cancellation $2M$]\label{ex:cv-fm}
For $M\ge1$ let $z\in[0,\frac34]$ be retained, $y\in[0,1]^2$ private, and
$\phi_M(z,y)=M(y_1+y_2-z)^2+(y_1-2z+\frac12)^2$. The retained-only part is
$Mz^2+(2z-\frac12)^2$, so $(H_0)_{zz}=2M+8$, and
$C=\left(\begin{smallmatrix}2M+2&2M\\2M&2M\end{smallmatrix}\right)\succ0$,
$\Gamma=(-2M-4,-2M)^{\T}$, $c=(1,0)^{\T}$. The unique responses and values
are
\[
\begin{array}{llll}
z\in[0,\tfrac14]:& y=(0,z),& \min_y\phi_M=(\tfrac12-2z)^2,& B^{\T}CB=2M,\\
z\in[\tfrac14,\tfrac12]:& y=(2z-\tfrac12,\tfrac12-z),& \min_y\phi_M=0,& B^{\T}CB=2M+8,\\
z\in[\tfrac12,\tfrac34]:& y=(\tfrac{(M+2)z-1/2}{M+1},0),& \min_y\phi_M=\tfrac{M(z-1/2)^2}{M+1},
& B^{\T}CB=\tfrac{2(M+2)^2}{M+1}.
\end{array}
\]
On the three intervals the private gradients are $(1-4z,0)$, $(0,0)$ and
$(0,2M(z-\frac12)/(M+1))$, all with the signs required at the active lower
bounds; in the third interval $y_1\in[\frac12,1)$. These are certificates.
By Proposition~\ref{prop:vf-curv}, $\sigma=2M$, and the certified curvature
is $2M+8-2M=8$, while the direct bound is $2M+8$. The piece second
derivatives $8,0,2M/(M+1)$ agree with $(H_0)_{zz}-B^{\T}CB$. The response is
unique and not affine, so no affine selector exists.
\end{example}

\begin{proposition}[A ladder with growing negative inertia]\label{prop:cv-ladder}
For $m\ge1$, $M\ge1$ and $a=\frac15$, let $z_i\in[0,\frac34]$,
$w_i\in[0,1]$ be retained, $y_{1i},y_{2i}\in[0,1]$ private, and
\[
F=\sum_{i=1}^m\Bigl[z_i^2+\phi_M(z_i,y_{1i},y_{2i})-w_i^2+3w_i-(z_i-a)w_i\Bigr]
+\frac1{16}\sum_{i=1}^{m-1}\Bigl[(z_{i+1}-z_i)^2+(w_{i+1}-w_i)^2\Bigr].
\]
Then $F$ has the unique minimizer $z_i=a$, $w_i=0$, $y_{1i}=0$,
$y_{2i}=a$, with $\OPT=m/20$; growth holds with constant $\frac1{12}$, and
no growth constant exceeds $\frac{33}{16}$. Its Hessian has exactly $m$
negative eigenvalues and $2\le\nu\le\sqrt5$. With the private blocks
$(y_{1i},y_{2i})$, the retained bags $\{z_i,w_i,z_{i+1},w_{i+1}\}$ have size
at most four, and the certificates of Example~\ref{ex:cv-fm} give
$L^+\le\frac{41}4$. The direct diagonal at $z_i$ is at least $2M+10$.
\end{proposition}

\begin{proof}
For one index put $t=z-a$, $\alpha_1=y_1$, $\alpha_2=y_2-a$,
$\rho_1=\alpha_1+\alpha_2-t$ and $\rho_2=\alpha_1-2t$. Direct expansion
gives
\[
z^2+\phi_M-\tfrac1{20}=t^2+M\rho_1^2+\rho_2^2+\tfrac{\alpha_1}5.
\]
Also $-w^2+3w-tw\ge w^2$, because $w(3-2w-t)\ge0$ for $w\in[0,1]$ and
$t\le\frac{11}{20}$. Since $\alpha_1\ge0$ and the path terms are
nonnegative,
\[
F-\tfrac m{20}\ge\sum_i\bigl(t_i^2+M\rho_{1i}^2+\rho_{2i}^2+w_i^2\bigr).
\]
From $\alpha_1=\rho_2+2t$ and $\alpha_2=\rho_1-\rho_2-t$,
$\alpha_1^2\le2\rho_2^2+8t^2$ and $\alpha_2^2\le3(\rho_1^2+\rho_2^2+t^2)$, so
$t^2+\alpha_1^2+\alpha_2^2\le12(t^2+\rho_1^2+\rho_2^2)\le12(t^2+M\rho_1^2+\rho_2^2)$.
Hence $F-\OPT\ge\frac1{12}\norm{x-x^*}^2$, and $F=\OPT$ forces
$t=\rho_1=\rho_2=w=0$ and $\alpha_1=\alpha_2=0$.

The terms $z_i^2-w_i^2-(z_i-a)w_i$ contribute the block
$\left(\begin{smallmatrix}2&-1\\-1&-2\end{smallmatrix}\right)$ on each pair
$(z_i,w_i)$, with eigenvalues $\pm\sqrt5$. All other quadratic terms, namely
the squares in $\phi_M$ and the path terms, are positive semidefinite. So
$\lambda_{\min}\ge-\sqrt5$. The direction with $w=m^{-1/2}(1,\ldots,1)$ and
all other coordinates zero has Rayleigh quotient $-2$, so $\nu\ge2$. On the
$m$-dimensional $w$-subspace the form is $-2I+\frac18L_{\mathrm{path}}$, where
$L_{\mathrm{path}}$ is the Laplacian of the path, and
$\frac18\norm{L_{\mathrm{path}}}\le\frac12$, so it is negative definite; on
the subspace $\{w=0\}$, of codimension $m$, the form is a sum of squares.
Hence the negative inertia is exactly $m$.

The retained diagonal at $z_i$ is
$(H_0)_{z_iz_i}=2+(2M+8)+\frac18\deg_i$ with $\deg_i\le2$, and at $w_i$ it
is $-2+\frac18\deg_i<0$. Each block has $\sigma=2M$ by
Example~\ref{ex:cv-fm}, so $L_{z_i}\le10+\frac14=\frac{41}4$ and
$L_{w_i}<0$. Raising $w_1$ from $0$ to $1$ at the minimizer increases $F$ by
at most $2+\frac1{16}$, which shows that no growth constant exceeds
$\frac{33}{16}$.
\end{proof}

Thus the certified ratio is $L^+/g\le123$ for all $m$ and $M$, whereas the
direct coordinate ratio is at least $(2M+10)\cdot16/33$. Every instance also
has the explicit nonnegative expansion above; the family separates
parameters and is not offered as a hard benchmark.

\begin{remark}[The certified cancellation is unstable]\label{rem:cv-unstable}
For the block objective $M(y-w)^2+\epsilon y$ with $y,w\in[0,1]$, the
response is $y=\max\{0,w-\epsilon/(2M)\}$. For $\epsilon=0$ it is $y=w$, so
$\sigma=2M$. For every $\epsilon>0$ the piece $[0,\epsilon/(2M)]$ has
response $y=0$, so $\sigma=0$ by Proposition~\ref{prop:vf-curv}(c), and the
value has second derivative $2M$ there. The full Hessian, and hence $\nu$,
does not depend on $\epsilon$. The same happens near $z=0$ in
Example~\ref{ex:cv-fm} if $\epsilon y_2$ is added: on a short interval both
private coordinates sit at their lower bounds, and $\sigma=0$. The certified
curvature is a maximum over pieces, however small the piece.
\end{remark}

\subsection{Proof of Proposition~\ref{prop:cv-limit}}

Put $\xi_1=x-1\in[-1,1]$ and $\xi_2=y-1\in[0,2]$, so
$F_M=M(\xi_1-\xi_2)^2-\frac18(\xi_1+\xi_2)^2+2\xi_2$. Using $M\ge1$ and
$2\xi_2\ge\xi_2^2$ on $[0,2]$,
\[
F_M-\tfrac18(\xi_1^2+\xi_2^2)\ge\tfrac34\xi_1^2-\tfrac94\xi_1\xi_2+\tfrac74\xi_2^2
=\tfrac34\bigl(\xi_1-\tfrac32\xi_2\bigr)^2+\tfrac1{16}\xi_2^2\ge0.
\]
This proves growth $\frac18$, uniqueness and $\OPT=0$. The Hessian
$\left(\begin{smallmatrix}2M-1/4&-2M-1/4\\-2M-1/4&2M-1/4\end{smallmatrix}\right)$
has eigenvalues $4M$ on $(1,-1)$ and $-\frac12$ on $(1,1)$; so $\nu=\frac12$,
and $\mathcal K=\emptyset$ is not admissible. The admissible sets are
$\{x,y\}$, $\{y\}$ and $\{x\}$. Note $F_M(2,2)=\frac32$.

$\mathcal K=\{x,y\}$: $V=F_M$ has $L_{\mathcal K}=2M-\frac14$. Under
$x=\omega_1x'$, $y=\omega_2y'$ the coordinate curvatures become
$\omega_i^2(2M-\frac14)$, and the point $(2,2)$ shows that every growth
constant in the new coordinates is at most
$\frac32/(\omega_1^{-2}+\omega_2^{-2})\le\frac32\min_i\omega_i^2$.

$\mathcal K=\{y\}$: $F_M$ is strictly convex in $x$, and
$\partial_xF_M(2,y)=2M(2-y)-y/4\le0$ for $y\ge y_c=4M/(2M+\frac14)$, where
$1<y_c<2$. So on $[y_c,3]$ the response is $x=2$, and $V(y)=F_M(2,y)$ has
second derivative $2M-\frac14$. Hence $L_{\mathcal K}\ge2M-\frac14$, while
$V(2)\le\frac32$ and $V(1)=0$ give $g_{\mathcal K}\le\frac32$.

$\mathcal K=\{x\}$: $F_M$ is strictly convex in $y$, and
$\partial_yF_M(x,1)=2-(2M+\frac14)(x-1)\ge0$ for
$x\le x_c=1+2/(2M+\frac14)$, where $1<x_c\le2$. So on $[0,x_c]$ the response
is $y=1$, and $V(x)=(M-\frac18)(x-1)^2$ has second derivative $2M-\frac14$.
Again $V(2)\le\frac32$ gives $g_{\mathcal K}\le\frac32$.

In one dimension a rescaling multiplies curvature and growth by the same
factor. In all cases
$L_{\mathcal K}/g_{\mathcal K}\ge(2M-\frac14)/\frac32=(4M-\frac12)/3\ge M$.
\qed
